% =====================================================================
%  CARNET D'ENTRAÎNEMENT — FICHE 6 : Les limites
%  Thème 1 — EXERCICES — Niveau FACILE
% =====================================================================
\documentclass[11pt,a4paper]{article}
\usepackage[utf8]{inputenc}
\usepackage[T1]{fontenc}
\usepackage{lmodern}
\usepackage[french]{babel}
\usepackage{amsmath,amssymb,mathtools}
\usepackage{icomma}
\usepackage{xcolor}
\usepackage{tikz}
\usepackage{pgfplots}
\usepackage[most]{tcolorbox}
\usepackage{enumitem}
\usepackage{array}
\usepackage{booktabs}
\usepackage{fancyhdr}
\usepackage[hidelinks]{hyperref}
\usetikzlibrary{arrows.meta,calc,intersections}
\pgfplotsset{compat=1.18}
\usepackage[a4paper,margin=2.2cm,headheight=26pt,headsep=10pt]{geometry}

\definecolor{primary}{RGB}{223,87,99}
\definecolor{primarydark}{RGB}{150,42,55}
\definecolor{excol}{RGB}{0,135,90}
\definecolor{curvecol}{RGB}{40,80,190}

\tcbset{boxbase/.style={enhanced,breakable,boxrule=0.9pt,arc=2.5pt,
   left=3mm,right=3mm,top=2.3mm,bottom=2.3mm,
   fonttitle=\bfseries,coltitle=white,toptitle=0.6mm,bottomtitle=0.6mm}}
\newtcolorbox[auto counter]{exercice}[1][]{boxbase,colback=primary!4,colframe=primary,
  title={\textsc{Exercice}~\thetcbcounter\ifx\relax#1\relax\relax\else\textmd{ --- \itshape #1}\fi}}

\tikzset{
  axe/.style={->,>=Stealth,thick,black!75},
  courbe/.style={line width=1.1pt,curvecol},
  asy/.style={dashed,black!55,line width=0.8pt},
  proj/.style={dashed,black!55,line width=0.6pt},
  pt/.style={circle,fill=black,inner sep=1.3pt},
  ptouvert/.style={circle,draw,fill=white,inner sep=1.1pt,line width=0.7pt},
}

\pagestyle{fancy}
\fancyhf{}
\fancyhead[L]{\small\textcolor{primarydark}{\textbf{Fiche 6} — Les limites}}
\fancyhead[R]{\small\textcolor{primarydark}{\textsc{Thème 1 — Facile}}}
\fancyfoot[C]{\small\thepage}
\renewcommand{\headrulewidth}{0.8pt}
\renewcommand{\headrule}{\hbox to\headwidth{\color{primary}\leaders\hrule height \headrulewidth\hfill}}

\newcommand{\R}{\mathbb{R}}

\begin{document}

\begin{center}
{\Large\bfseries\color{primarydark}Thème 1 — Notations, lecture \& calcul direct}\\[2pt]
{\large\color{primary}Exercices — Niveau Facile}
\end{center}
\vspace{4pt}
{\small\itshape Objectif : lire une notation de limite et calculer une limite par remplacement direct
quand la fonction est définie. On travaille \textbf{une} mécanique à la fois.}
\vspace{6pt}

\begin{exercice}[calcul direct par remplacement]
Toutes ces fonctions sont définies au point visé. Calculer chaque limite en remplaçant $x$ par la cible.
\begin{enumerate}[label=\alph*),leftmargin=2em,itemsep=3pt]
  \item $\displaystyle\lim_{x\to 2}\bigl(3x^2-x+4\bigr)$
  \item $\displaystyle\lim_{x\to -1}\frac{2x+5}{x^2+3}$
  \item $\displaystyle\lim_{x\to 4}\bigl(\sqrt{x}+x\bigr)$
  \item $\displaystyle\lim_{x\to 0}\bigl(x^3+2x-7\bigr)$
\end{enumerate}
\end{exercice}

\begin{exercice}[lire et écrire une notation]
\begin{enumerate}[label=\alph*),leftmargin=2em,itemsep=3pt]
  \item Traduire en une phrase : $\displaystyle\lim_{x\to +\infty} f(x)=3$.
  \item Traduire en une phrase : $\displaystyle\lim_{x\to 2^{-}} f(x)=-\infty$.
  \item Écrire avec le symbole $\lim$ : « lorsque $x$ se rapproche de $5$ par la droite, $g(x)$ tend vers $0$ ».
\end{enumerate}
\end{exercice}

\begin{exercice}[limites à gauche et à droite d'une fonction par morceaux]
Soit $f$ définie par
\[ f(x)=\begin{cases} x+1 & \text{si } x<2,\\[2pt] 3x-2 & \text{si } x\geq 2.\end{cases} \]
\begin{enumerate}[label=\alph*),leftmargin=2em,itemsep=3pt]
  \item Calculer $\displaystyle\lim_{x\to 2^{-}} f(x)$ (on utilise le morceau valable pour $x<2$).
  \item Calculer $\displaystyle\lim_{x\to 2^{+}} f(x)$ (on utilise le morceau valable pour $x\geq 2$).
  \item Les deux résultats sont-ils égaux ? $f$ a-t-elle une limite (globale) en $2$ ?
\end{enumerate}
\end{exercice}

\begin{exercice}[lecture graphique]
La courbe de $f$ est donnée ci-dessous. Lire graphiquement les limites demandées.
\begin{center}
\begin{tikzpicture}[scale=0.9]
\begin{axis}[width=10cm,height=5.6cm,axis lines=middle,xlabel={$x$},ylabel={$y$},
   xmin=-5.5,xmax=5.6,ymin=-1.2,ymax=3.4,xtick={-4,-2,2,4},ytick={1,2,3},
   tick label style={font=\scriptsize}]
  % courbe : 2 - 2/(x^2+1)  -> tend vers 2 en +-infini, vaut 0 en 0
  \addplot[courbe,domain=-5.4:5.4,samples=120]{2-2/(x^2+1)};
  \draw[asy,primary](axis cs:-5.5,2)--(axis cs:5.6,2);
  \node[primary,font=\scriptsize,anchor=south west] at (axis cs:-5.3,2.05){$y=2$};
\end{axis}
\end{tikzpicture}
\end{center}
\begin{enumerate}[label=\alph*),leftmargin=2em,itemsep=3pt]
  \item $\displaystyle\lim_{x\to +\infty} f(x)=\ ?$ \qquad
  \item $\displaystyle\lim_{x\to -\infty} f(x)=\ ?$ \qquad
  \item $\displaystyle\lim_{x\to 0} f(x)=\ ?$
\end{enumerate}
\end{exercice}

\begin{exercice}[vérifier une continuité]
Soit $g(x)=\dfrac{x^2-1}{x+3}$.
\begin{enumerate}[label=\alph*),leftmargin=2em,itemsep=3pt]
  \item Le point $x_0=1$ appartient-il au domaine de $g$ ? (le dénominateur s'annule-t-il ?)
  \item Calculer $\displaystyle\lim_{x\to 1} g(x)$ et la valeur $g(1)$.
  \item En déduire si $g$ est continue en $1$ (comparer la limite et $g(1)$).
\end{enumerate}
\end{exercice}

\end{document}
